\section{Planar projection and product-pinning proofs}
\label{sec:projection-proofs}

This section proves the two planar statements used in
Section~\ref{sec:projections}. We first derive the robust projection
bound from the quantitative Furstenberg theorem. We then use that bound
to construct the pair law whose product slope is nonconcentrated.

\subsection{The robust projection bound}
\label{sec:robust-projection-proof}

We prove Lemma~\ref{lem:robust-projection}. Its parameters $s,t,u$
are fixed, and $h_{\mathrm{pr}}=\min\{t,(s+t)/2,1\}$ is the planar
projection exponent. A positive product hit probability would
produce a Furstenberg configuration with too few tubes: comparable
slope masses and incidence counts give the required quantitative
point and direction bounds.

\begin{proof}
Put $\Delta=h_{\mathrm{pr}}-u>0$ and $\varepsilon=\Delta/4$. Let
$\eta_{\mathrm{pr}}=\eta_{\mathrm{pr}}(\varepsilon,s,t)>0$ be the Frostman tolerance in
\citet[Theorem~4.1]{RW2025}. We may take
\begin{equation}\label{eq:robust-projection-kappa}
 \kappa_0=\min\{1/20,\Delta/16,\eta_{\mathrm{pr}}/12\}.
\end{equation}
All constants below may depend on $R,C_0,s,t$, but not on the
number of labels or on their weights. We decrease $\delta_0$
finitely many times.

First put the incidences in the bounded coordinates used for
dyadic tubes. Set $B_0=R+1$ and
\[
 v=\frac{y+B_0}{2B_0},\qquad
 w=\frac{x+B_0}{2B_0^2},\qquad
 a_\lambda=-\frac{\vartheta(\lambda)}{B_0}.
\]
Then $(v,w)\in(0,1)^2$, $a_\lambda\in(-1,1)$, and
\[
 w-a_\lambda v
 =\frac{x+\vartheta(\lambda)y+B_0(1+\vartheta(\lambda))}
        {2B_0^2}.
\]
Apply this affine transformation to $A_\lambda$, writing the
result as $J_\lambda$. Each $J_\lambda$ has at most
$C\delta^{-u}$ constituent intervals of length at most $C\delta$.
The transformed point and slope measures satisfy
\eqref{eq:robust-projection-frostman} with an additional fixed
factor $C$. We keep the notation $\mu$ for the transformed point
measure.

Choose a dyadic number $\varrho$ with
$\delta\le\varrho<2\delta$, and put
\[
 L=8\lceil\log_2(1/\delta)\rceil+8,
 \qquad q_*=\delta^{2\kappa}.
\]
Suppose that the left side of \eqref{eq:robust-projection-product}
is at least $q_*$. Write
$f(\lambda)=\mu\{(v,w):w-a_\lambda v\in J_\lambda\}$.
Partition the slopes into dyadic $\varrho$-intervals $I$. Give
$I$ its total label weight
$W_I=\nu\{\lambda:a_\lambda\in I\}$, and choose a label
$\lambda_I$ maximizing $f$ in that interval. There are finitely
many labels, so
\begin{equation}\label{eq:robust-projection-aggregate}
 \sum_I W_If(\lambda_I)\ge q_*.
\end{equation}
There are at most $C\delta^{-1}$ occupied slope intervals.
Discarding those with $W_I<\delta^4$ loses at most
$C\delta^3\le q_*/2$ from \eqref{eq:robust-projection-aggregate}.
The remaining weights have at most $L$ dyadic levels. One level,
consisting of $n$ slope intervals, satisfies
\[
 H_*:=\sum_{I\text{ in the level}}W_If(\lambda_I)
 \ge\frac{q_*}{2L}.
\]
Let $a_* =\sum_I W_I$ on this level. Its weights are within a
factor of two, and $H_*\le a_*\le1$. Consequently the uniform
distribution of its slope intervals satisfies
\begin{equation}\label{eq:robust-projection-slopes}
 \frac{\#\{I:I\cap B(b,r)\ne\varnothing\}}n
 \le CL\delta^{-3\kappa}r^s
 \quad(\varrho\le r\le1).
\end{equation}
Indeed, relative counting costs at most $2/a_*$ times the
original slope mass of a fixed enlargement of the ball, and
$a_*\ge q_*/(2L)$. For enlarged radii exceeding one, the same
bound follows from the total mass bound.

Let $d(p)$ be the number of selected representative tests
satisfied by $p=(v,w)$. Comparability of the slope weights gives
\[
 \int\frac{d(p)}n\,d\mu(p)
 \ge\frac{H_*}{2a_*}\ge\frac{q_*}{4L}.
\]
Since $0\le d(p)/n\le1$ and $\mu(\R^2)\le1$, the set
$G_0=\{p:d(p)\ge q_*n/(8L)\}$ has mass at least $q_*/(8L)$.
Pigeonhole its positive integer degrees into at most $L$
dyadic levels. This gives a Borel set $G\subset G_0$ and an
integer $M\ge1$ such that
\begin{equation}\label{eq:robust-projection-degrees}
 \mu(G)\ge\frac{q_*}{8L^2},\qquad
 M\le d(p)<2M\ (p\in G),\qquad
 M\ge\frac{q_*n}{16L}.
\end{equation}

Partition the point square into dyadic $\varrho$-squares and give
$Q$ the weight $\mu(G\cap Q)$. There are at most
$C\delta^{-2}$ squares. Discard weights below $\delta^4$; their
total is at most $C\delta^2\le q_*/(16L^2)$. A dyadic weight
pigeonhole gives a family $\mathcal P$ of squares with comparable
weights and total retained mass
\begin{equation}\label{eq:robust-projection-point-mass}
 b_*\ge\frac{q_*}{16L^3}.
\end{equation}
For each $Q\in\mathcal P$ choose $p_Q\in G\cap Q$.
Comparison between counting and these comparable cell weights
shows that $\mathcal P$ has relative $t$-Frostman constant
\begin{equation}\label{eq:robust-projection-point-constant}
 CL^3\delta^{-3\kappa}
\end{equation}
down to $\varrho$: passage to relative counts costs at most
$2/b_*$, and the original measure bounds the mass of a fixed
enlargement of each testing ball.

Here a dyadic $\varrho$-tube is the union of lines
$w=av+b$ with $(a,b)$ in a fixed dyadic parameter square of side
$\varrho$. For each selected slope representative $a_I$ and each
$b'\in J_{\lambda_I}$, include the tube whose parameter square
contains $(a_I,b')$. Call the resulting family $\mathcal T$.
Every $J_{\lambda_I}$ meets at most $C\delta^{-u}$ intercept
intervals, so
\begin{equation}\label{eq:robust-projection-tube-upper}
 |\mathcal T|\le Cn\delta^{-u}.
\end{equation}
For $Q\in\mathcal P$ and each representative test satisfied by
$p_Q=(v_Q,w_Q)$, select the tube with parameter cell containing
$(a_I,w_Q-a_Iv_Q)$. This tube intersects $Q$. Different slope
cells give different tubes. Denote this fiber by $\mathcal T(Q)$;
then $M\le|\mathcal T(Q)|<2M$.

A ball of radius $r$ in tube-parameter space restricts its slope
coordinate to an interval of comparable radius. By
\eqref{eq:robust-projection-slopes} and
\eqref{eq:robust-projection-degrees},
\begin{equation}\label{eq:robust-projection-fiber-constant}
 \frac{\#\{\text{parameter cells of }\mathcal T(Q)
                     \text{ meeting }B(z,r)\}}
      {|\mathcal T(Q)|}
 \le CL^2\delta^{-5\kappa}r^s
 \quad(\varrho\le r\le1).
\end{equation}
Thus the base cells and selected incident fibers form exactly
the quantitative configuration of
\citet[Definition~3.1 and Theorem~4.1]{RW2025}. The selected fiber
cardinalities lie between $M$ and $2M$, as
permitted by the cited theorem. For sufficiently small $\delta$, the
constants in \eqref{eq:robust-projection-point-constant} and
\eqref{eq:robust-projection-fiber-constant} are at most
$\varrho^{-\eta_{\mathrm{pr}}}$. In detail, arrange $CL^3\le\delta^{-\kappa}$
and use $6\kappa\le\eta_{\mathrm{pr}}/2$ together with
$\varrho<2\delta$.

The cited theorem now gives
\[
 |\mathcal T|\ge c_\varepsilon M\varrho^{-h_{\mathrm{pr}}+\varepsilon}
 \ge\frac{c'_\varepsilon}{16L}
             n\delta^{-h_{\mathrm{pr}}+\varepsilon+2\kappa}.
\]
Combining this with \eqref{eq:robust-projection-tube-upper} yields
\[
 1\ge\frac{c''_\varepsilon}{L}
             \delta^{-(\Delta-\varepsilon-2\kappa)}.
\]
But \eqref{eq:robust-projection-kappa} implies
$\Delta-\varepsilon-2\kappa\ge5\Delta/8>0$, so the right
side tends to infinity. This proves
\eqref{eq:robust-projection-product}.
Finally, \eqref{eq:robust-projection-subset} would make the product
hit probability at least $\nu(B)\delta^\kappa\ge\delta^{2\kappa}$,
a contradiction.
\end{proof}

\begin{remark}[Measures, labels, and fixed charts]
\label{rem:robust-projection-extensions}
The point measure need not be atomic: finite point families
appear only after restricted cell masses are pigeonholed. The
statement also holds for a general measurable label space if
the slope map and the target incidence set in the product of
label and point spaces are jointly measurable. In each of the
finitely many slope cells, choose a label whose hit probability
is at least half the conditional average. This replaces
\eqref{eq:robust-projection-aggregate} by the same inequality
with a factor $1/2$, absorbed by the constants in the proof.

A fixed finite collection of bounded projection charts is
permitted: a violating product hit probability must have a
fixed fraction in one chart, and restrictions are still
subprobability measures. Fixed affine coordinate changes,
label-dependent translations of the target, and scalar
normalizations bounded above and away from zero also only
alter fixed constants and the small-scale threshold.

At a local scale $\delta=(\delta')^{h_1}$, an inherited bound
$(\delta')^{-a}r^t$ fits
\eqref{eq:robust-projection-frostman} when $a\le\kappa h_1$.
This explicitly determines how small losses must be chosen
on each fixed finite depth ladder. If a measure to be tested
is bounded above by $\delta^{-b}$ times the independent product,
where $b<2\kappa$, its hit probability is at most
$\delta^{2\kappa-b}$. No independence of a preceding coupled
pair measure is asserted.
\end{remark}


\subsection{Constructing the product-pinning law}
\label{sec:product-pinning-proof}

We now prove Lemma~\ref{lem:pinning}. There are two possibilities.
If a thin line neighborhood carries substantial mass, the coordinate
Frostman bounds keep its line away from the coordinate directions;
on that cluster the product map is a nondegenerate quadratic in time.
Otherwise, avoidance of line neighborhoods gives an initial positive
ray exponent. Repeated planar projection estimates improve that
exponent toward the dimension of the parameter measure.

This ray-improvement argument adapts the projection bootstrap of
\citet{SW2025}, developed further by \citet{OSW2024}. We prove the finite-scale
form required here, including the product map and bounds for both
single-event marginals.

We use unoriented directions, with their usual angular metric, and
write
\[
 R_\xi(\upsilon)=[\upsilon-\xi]\in\mathbb P^1(\R)
 \qquad(\upsilon\ne\xi).
\]
The next lemma supplies the nonlinear part of the proof.
Absolute section masses, rather than normalized conditional
probabilities, make the symmetry and deletion estimates transparent.
The ray improvement and ensuing pinning iteration are inspired by the
thin-tube bootstrap of Orponen--Shmerkin--Wang
\citep[Section~1.3, Lemma~2.8 and Corollary~2.18]{OSW2024}.
The regular finite-scale and product-pinning arguments are proved here,
using the separately credited Ren--Wang projection input; these
interfaces are not invoked as consequences of the cited bootstrap.

\begin{lemma}[Improving ray sections]
\label{lem:pr-ray-improvement}
Let $0<s\le1$ and $\delta\to0$. For each $\delta$, let $\mu$ be a
finitely supported probability on a bounded square
such that every occupied square in a fixed nested dyadic grid
of side comparable to $\delta^u$ has mass
$\delta^{su+o(1)}$, uniformly for $0\le u\le1$.
Suppose $G\subset\supp\mu\times\supp\mu$ is symmetric,
$\mu^2(G)=1-o(1)$, and, for some $0<\sigma<s$,
\begin{equation}
 \mu\{\upsilon:(\xi,\upsilon)\in G,
             R_\xi(\upsilon)\in B_r\}
 \le\delta^{-o(1)}r^\sigma
 \qquad(\delta\le r\le1)
\label{pr:old-ray-bound}
\end{equation}
for every anchor $\xi$ and every direction ball $B_r$.
For every
\[
 0<\sigma_*<\min\{s,(s+\sigma)/2\},
\]
one can restrict to a symmetric relation $G_*\subset G$ of
probability $1-o(1)$ on which \eqref{pr:old-ray-bound} holds with
$\sigma_*$ in place of $\sigma$.
One can instead obtain the same absolute section bound for
$\mathfrak p_\xi(\upsilon)$, with intervals in place of direction
balls. All assertions are uniform for ensembles satisfying the
hypotheses uniformly.
\end{lemma}

\begin{proof}
We localize the second point to a small square and linearize the
output map there. Symmetry supplies the reverse rays from the second
point to the anchors; those rays provide the directions for the planar
projection estimate. A joint pruning then makes the local estimates
iterable over a finite ladder of depths.

Choose
\begin{equation}
 \sigma_*<\sigma_m<\min\{s,(s+\sigma)/2\}.
\label{pr:intermediate-exponent}
\end{equation}
We first prove the bound at a fixed depth $T\in(0,1]$.
Choose a small starting depth $d_0>0$ and a finite ladder
\begin{equation}
 d_0<d_1<\cdots<d_n=T,
 \qquad d_{j+1}<2d_j.
\label{pr:ladder}
\end{equation}
The choices of $d_0$ and of a further small parameter $\lambda>0$
will be specified below. Delete pairs of distance less than
$\delta^\lambda$. Their probability is at most
$\delta^{s\lambda-o(1)}$, by the square-mass upper bounds.

\paragraph{A local projection estimate.}
Consider one step $d_1\to d_2$ of the ladder, and put
$h_1=d_2-d_1$. Choose $\lambda$ small enough that
\begin{equation}
 d_1-\lambda>h_1,
 \qquad 2d_1-2\lambda>d_2.
\label{pr:geometric-margins}
\end{equation}
Let $J$ be an occupied square of side $\delta^{d_1}$ and let
$\mu_J=\mu|_J/\mu(J)$, rescaled to a bounded unit square.
For $0\le u\le h_1$, the ratio of child and parent masses is
$\delta^{su+o(1)}$. A bounded square covering therefore makes
$\mu_J$ an $s$-Frostman probability through resolution
\[
 \delta_J=\delta^{h_1}.
\]
This uses the lower bound on $\mu(J)$ as well as the upper bound on
its children.

For a fixed $\upsilon\in J$, consider the anchors $\xi$ paired
with $\upsilon$ by the old relation and satisfying the separation
condition. Symmetry and \eqref{pr:old-ray-bound} give an absolute
$\sigma$-Frostman bound on their reverse rays.
Replacing $\upsilon$ by the center of $J$ changes the direction by
$O(\delta^{d_1-\lambda})$. By
\eqref{pr:geometric-margins}, this is smaller than $\delta_J$.
The same absolute bound thus holds for the directions perpendicular
to the rays from these anchors to the center of $J$.

Apply Lemma~\ref{lem:robust-projection} at resolution $\delta_J$,
with point exponent $s$, direction exponent $\sigma$, and target
exponent $\sigma_m$. Fix its positive power allowance $\kappa$.
A reverse section of mass less than $\delta_J^{\kappa/4}$ may be
discarded at total pair cost at most $\delta_J^{\kappa/4}$.
For larger sections, normalization increases the Frostman constant
by at most $\delta_J^{-\kappa/4}$. All inherited subpower constants
are smaller than $\delta_J^{-\kappa/4}$ for sufficiently small
$\delta$, since $h_1$ is fixed. Decreasing $\kappa$ if necessary
leaves the loss required by the projection lemma.

For each direction, call a projection bin of width $\delta_J$
heavy if its $\mu_J$-mass exceeds $\delta_J^{\sigma_m}$.
The heavy bins number at most $O(\delta_J^{-\sigma_m})$.
The robust projection lemma implies that the directions for which
their union has mass at least $\delta_J^\kappa$ occupy at most a
fixed positive power of $\delta_J$ in every normalized reverse
section. One can use the product bound of that lemma and Markov's
inequality; shrinking $\kappa$ accommodates constants.
Call these directions exceptional for $J$.

There are two different removals here. First remove pairs whose
anchor direction is exceptional for the cell containing the other
point. For each fixed $\upsilon$, its reverse-anchor mass is
power-small; integration against $\mu(d\upsilon)$ gives a
power-small total loss. Second, for each nonexceptional $(\xi,J)$,
remove the points in $J$ lying in heavy projection bins.
Their \emph{unconditional} $\mu$-mass is at most
$\delta_J^\kappa\mu(J)$. Integrating over $\xi$ and summing over
the disjoint cells $J$ again gives a power-small loss.
Exceptionality is a test of $J$ and the anchor direction, independent
of the point within $J$. Neither removal assumes independence on
the old relation $G$.

For the ray map, its derivative on $J$ is a perpendicular scalar
projection multiplied by the inverse distance to the anchor.
That multiplier is bounded below by a positive constant, since the
ambient square is bounded. Its quadratic error is
$O(\delta^{2d_1-2\lambda})$. Consequently
\eqref{pr:geometric-margins} implies the following local bound on
the surviving section:
\begin{equation}
 \mu\{\upsilon\in J:(\xi,\upsilon)\text{ survives},
                R_\xi(\upsilon)\in B_{\delta^{d_2}}\}
 \le C\delta^{\sigma_m h_1}\mu(J).
\label{pr:local-ray}
\end{equation}
Finite angular charts and overlapping projection grids only change
$C$.

For the product map, the derivative is
\begin{equation}
 \nabla_\upsilon\mathfrak p_\xi(\upsilon)
   =(\vartheta_\upsilon-\vartheta_\xi,
                         t_\upsilon-t_\xi).
\label{pr:product-gradient}
\end{equation}
Its direction is the swapped ray, whose direction law has the same
Frostman exponent. Thus the product test still uses the old ray bound
as its direction input; it does not require a previous product bound.
Its size is at least $\delta^\lambda$.
An output interval of width $\delta^{d_2}$ consequently uses at
most $\delta^{-O(\lambda)}$ projection bins at resolution
$\delta_J$. The quadratic error is $O(\delta^{2d_1})$.
Thus the product version of \eqref{pr:local-ray} has right side
$C\delta^{\sigma_mh_1-C\lambda}\mu(J)$.
In what follows this slightly weaker bound works for both maps.

\paragraph{Pruning and iteration.}
Perform these removals for every step in the finite union of tested
ladders, also with the legs interchanged. After all these removals,
perform one joint pruning process over every tested level and both
legs: if, for an anchor and a tested
cell $J$, the remaining nonempty section has mass less than
$\delta^\lambda\mu(J)$, delete it. Continue with either leg until
no deficient section remains at any tested level. Each
anchor--cell combination is charged at most once,
because it is empty after its deletion. Integrating over anchors
and summing over $J$ bounds the charge by $\delta^\lambda$ per
tested level and per leg, even if the deletions cause a cascade.
There are finitely many levels, so the total loss tends to zero.
All operations are finite on the finitely supported measures. Any
subrelation satisfying every section threshold survives each deletion,
so the terminal relation is the unique largest subrelation with those
thresholds. Interchanging the legs preserves both the initial relation
and the tests; hence the terminal relation is symmetric.

For the final relation, let $Q(d)$ be the largest absolute section
mass in an output ball of radius $\delta^d$. The retained mass in a
nonempty anchor section of $J$ is now at
least $\delta^\lambda\mu(J)$. The local estimate divided by this
lower bound is a factor
\[
 C\delta^{\sigma_mh_1-C_1\lambda}
\]
times that section mass. If the section has an output in a chosen
fine ball, all its outputs lie in an enlarged parent ball of radius
$O(\delta^{d_1-\lambda})$. That ball is coverable by
$\delta^{-O(\lambda)}$ balls of radius $\delta^{d_1}$.
The cells are disjoint, so summing their section masses gives the
recurrence
\begin{equation}
 Q(d_2)\le C\delta^{\sigma_m(d_2-d_1)-C_2\lambda}Q(d_1).
\label{pr:ladder-recurrence}
\end{equation}
Restrictions made at other levels preserve every upper bound.

Starting with $Q(d_0)\le1$, the ladder gives
\[
 Q(T)\le C^n
       \delta^{\sigma_m(T-d_0)-C_2n\lambda-o(1)}.
\]
First choose $d_0$ so that
$\sigma_m d_0<(\sigma_m-\sigma_*)T/4$.
After fixing the ladder, choose $\lambda$ so that
$C_2n\lambda<(\sigma_m-\sigma_*)T/4$ and all the strict
inequalities \eqref{pr:geometric-margins} hold.
The positive remaining margin absorbs $C^n$ and the subpower
errors. This proves $Q(T)\le\delta^{\sigma_*T}$ for small
$\delta$.

The same construction works for finitely many target depths at once:
take the union of their ladders and choose all loss parameters below
the minimum of their finitely many positive allowances.
To obtain all scales, take successively finer finite rational meshes
of target depths. For each fixed mesh first make $\delta$ small
enough that its finite collection of projection and deletion bounds
holds and its total deleted mass is as small as prescribed.
Then let the mesh refine sufficiently slowly. Between tested
depths monotonicity costs at most $\delta^{-\sigma_*e}$, where
$e$ is the mesh spacing; at depths below the first mesh point the
trivial mass bound has the same cost. Since $e\to0$, this gives the
claimed subpower constant at all radii. This order never compares
an uncontrolled error with a shrinking ladder step.
All bounds used above are uniform for the stipulated ensembles.
\end{proof}

\begin{proof}[Proof of Lemma~\ref{lem:pinning}]
It is enough to prove the assertion for small $\epsilon<s/10$;
we reserve a fixed fraction of $\epsilon$ for each normalization.
Choose a small constant $\eta_0>0$, depending on
$s,\gamma',\epsilon$, to be fixed in the first case below.

\paragraph{A line cluster.}
Suppose first that, for some fixed rational
$\zeta\in(0,\zeta_{\max}]$, a set of initial indices of probability
bounded below has a line neighborhood of width
$\delta=D^{-\zeta}$ carrying mass at least $\delta^{\eta_0}$.
Pass to that subsequence, restrict the initial index there, and
choose one such cluster $C_i$ on each index.
The normalization of the initial-index law costs a fixed factor.

Let $(v_t,v_\vartheta)$ be a unit vector parallel to its line.
The time and normal-coordinate ranges of the cluster have lengths
at most $C(|v_t|+\delta)$ and $C(|v_\vartheta|+\delta)$.
The coordinate bounds and its mass lower bound imply
\[
 |v_t|\ge\delta^{C\eta_0},\qquad
 |v_\vartheta|\ge\delta^{C\eta_0}
\]
for small $\delta$, with $C$ depending only on the fixed exponents.
Thus its slope $k_i=v_\vartheta/v_t$ satisfies
\begin{equation}
 \delta^{C\eta_0}\le |k_i|\le\delta^{-C\eta_0},
 \qquad
 \vartheta=k_i t+b_i+O(\delta^{1-C\eta_0})
 \quad\hbox{on }C_i.
\label{pr:cluster-slope}
\end{equation}
Here the subpower constants have been absorbed by slightly enlarging
$C$. Choose $K$ with $Ks>3$, and require $K\eta_0<1$ in the standing
choice of $\eta_0$. Sample two points from the normalized
cluster, restricting to
\[
 |t_\upsilon-t_\xi|\ge\delta^{K\eta_0}.
\]
For every anchor the rejected conditional mass is at most
$\delta^{(Ks-1)\eta_0-o(1)}=o(1)$.
The pair restriction can be normalized symmetrically.
Both marginals are bounded by
$(1-o(1))^{-1}\mu_i/\mu_i(C_i)$.

On these pairs,
\[
 \mathfrak p_\xi(\upsilon)
  =k_i(t_\upsilon-t_\xi)^2
       +O(\delta^{1-C\eta_0}).
\]
The quadratic is monotone on each separated branch and its
derivative there has magnitude at least $\delta^{C'\eta_0}$.
For every $r\ge\delta$, the preimage of an output interval of
radius $r$ is contained in at most two time intervals of total
length
\[
 C(r+\delta^{1-C\eta_0})\delta^{-C'\eta_0}
 \le Cr\delta^{-C''\eta_0}.
\]
Using the time Frostman bound before normalizing the cluster proves
\begin{equation}
 \Prob\{\mathfrak p_\xi(\upsilon)\in J_r\mid i,\xi\}
   \le C\delta^{-C'''\eta_0-o(1)}r^s
   \qquad(\delta\le r\le1).
\label{pr:cluster-final}
\end{equation}
The estimate includes $r=\delta$: the larger approximation error
costs the displayed small power, which has not been suppressed.
Fix $\eta_0$ sufficiently small compared with $\epsilon$ and all
the constants above, and then take $D$ large.
Equation~\eqref{pr:cluster-final} implies
\eqref{pr:pinned-frostman}; the two marginal losses, including the
initial-index restriction, are at most $D^{\epsilon\zeta}$.

\paragraph{No line cluster.}
In the other case, the proportion of indices with such a cluster
tends to zero at every fixed rational depth. Testing finite depth
meshes and discarding their exceptional indices, and only then
refining the meshes, gives uniformly on the retained indices
\begin{equation}
 \mu_i(\{\dist(\cdot,L)\le r\})
     \le D^{o(1)}r^{\eta_0}
 \quad(D^{-\zeta_{\max}}\le r\le1)
\label{pr:no-line-cluster}
\end{equation}
for all lines $L$. To verify the interpolation, replace a radius by
the next larger tested radius; if the depth mesh has spacing $e$,
the loss is at most $D^{\eta_0e}$. The discarded index probability
can tend to zero while $e\to0$, since every fixed finite test has
exceptional probability tending to zero.

Fix any rational $0<\zeta\le\zeta_{\max}$ and put
$\delta=D^{-\zeta}$. A ray-direction ball through an anchor is
contained in a line neighborhood of comparable width.
Equation~\eqref{pr:no-line-cluster} therefore supplies
\eqref{pr:old-ray-bound} with a positive exponent, decreasing
$\eta_0$ if necessary to make it smaller than $s$.
Start with all distinct pairs. Their probability tends to one,
because the mass of any point is at most that of its square at
resolution $\delta$, which is $\delta^{s-o(1)}$.

Repeatedly apply Lemma~\ref{lem:pr-ray-improvement}.
For example, choosing the next exponent halfway between the old
one and $\min\{s,(s+\sigma)/2\}$ decreases its distance to $s$
by a fixed factor. Finitely many iterations give a ray exponent
as close to $s$ as desired. Apply the product version once more to
obtain an absolute product-section exponent larger than
$s-\epsilon/2$, with subpower constant, on a symmetric relation of
probability $1-o(1)$.

Restrict first to anchors whose surviving section has mass at least
$1/2$, and then normalize on their sections. The discarded anchor
mass tends to zero. The first marginal is bounded by a constant
times $\mu_i$. The second marginal is also so bounded: its density
is an integral of section densities bounded by $2$, against a
first-anchor law bounded by a constant times $\mu_i$.
The product-section bound survives these normalizations.
Uniformity makes the same construction possible on every retained
initial index; the initial-index normalization is $1+o(1)$.
All fixed and subpower constants fit inside the reserved
$D^{\epsilon\zeta}$ allowance for sufficiently large $D$.
This proves the lemma and its independence from subsequent offset
tests.
\end{proof}


The two auxiliary inputs are now proved. Together with the same-index
information argument, they establish all the constraints in
Proposition~\ref{prop:projection-rates}. We turn to their incompatibility
with the stationary entropy demand.
